% Compatibilidad con instalaciones MiKTeX cuyo formato base es anterior a
% los paquetes Beamer/Hyperref instalados. En formatos recientes, estas
% definiciones no alteran nada porque \providecommand conserva las originales.
\providecommand{\IfFormatAtLeastTF}[3]{#3}
\providecommand{\IfFormatAtLeastT}[2]{}
\providecommand{\IfFormatAtLeastF}[2]{#2}
\providecommand{\IfDocumentMetadataTF}[2]{#2}
\providecommand{\IfDocumentMetadataT}[1]{}
\providecommand{\IfDocumentMetadataF}[1]{#1}
% Esta sesión usa únicamente PNG; evita cargar el conversor EPS de una
% instalación MiKTeX incompleta (grfext.sty no es necesario aquí).
\makeatletter
\@namedef{ver@epstopdf-base.sty}{9999/12/31}
\makeatother

\documentclass[aspectratio=169,11pt]{beamer}

% =============================================================================
% MACROECONOMIA II | LICENCIATURA EN ECONOMIA | CIDE
% SESION 3: EQUILIBRIO GENERAL EN UNA ECONOMIA CERRADA
% Compilacion normal: pdflatex (o latexmk).
% =============================================================================


% --- Idioma, tipografia y matematicas -----------------------------------------
\usepackage[utf8]{inputenc}
\IfFileExists{spanish.ldf}{
  \usepackage[spanish,es-nodecimaldot]{babel}
}{
  \usepackage[english]{babel}
}
% Traducciones de los entornos de Beamer
\uselanguage{Spanish}
\languagepath{Spanish}

\usepackage{graphicx}
\usepackage{amsmath,amssymb}
\IfFileExists{booktabs.sty}{
  \usepackage{booktabs}
}{
  \providecommand{\toprule}{\hline}
  \providecommand{\midrule}{\hline}
  \providecommand{\bottomrule}{\hline}
}
\providecommand{\addlinespace}[1][0.5em]{}
\usepackage{xcolor}
\usepackage{ragged2e}
\usepackage{etoolbox}


\newtheorem{proposition}[theorem]{\translate{Proposition}}
\newtheorem{remark}[theorem]{\translate{Remark}}
\newtheorem{assumption}[theorem]{Supuesto}

% --- Justificación del contenido ---------------------------------------------
% Beamer restablece \raggedright al iniciar listas, columnas y bloques. Estos
% parches sustituyen esa instrucción dentro de los comandos internos de Beamer.
\AtBeginDocument{\justifying}
\patchcmd{\itemize}{\raggedright}{\justifying}{}{}
\patchcmd{\enumerate}{\raggedright}{\justifying}{}{}
\patchcmd{\description}{\raggedright}{\justifying}{}{}

\addtobeamertemplate{block begin}{}{\justifying}
\addtobeamertemplate{block alerted begin}{}{\justifying}
\addtobeamertemplate{block example begin}{}{\justifying}

\makeatletter
\patchcmd{\beamer@columnenv}{\raggedright}{\justifying}{}{}
\patchcmd{\beamer@columncom}{\raggedright}{\justifying}{}{}
\makeatother

% --- Bibliografia opcional ----------------------------------------------------
\newif\ifusarBibliografia
\usarBibliografiafalse

\ifusarBibliografia
  \IfFileExists{csquotes.sty}{\usepackage{csquotes}}{}
  \IfFileExists{apa.bbx}{
    \usepackage[backend=biber,style=apa,sorting=nyt]{biblatex}
  }{
    \usepackage[backend=biber,style=authoryear,sorting=nyt]{biblatex}
  }
  \IfFileExists{References.bib}{
    \addbibresource{References.bib}
  }{
    \addbibresource{References.bib.bib}
  }
\fi

% --- Paleta: azul marino, blanco y grises neutros ----------------------------
\definecolor{CIDENavy}{RGB}{3,25,54}
\definecolor{CIDEBlue}{RGB}{5,35,75}
\definecolor{CIDEInk}{RGB}{25,34,45}
\definecolor{CIDEGray}{RGB}{101,112,126}
\definecolor{CIDEGrayLight}{RGB}{243,244,246}

\setbeamercolor{normal text}{fg=CIDEInk,bg=white}
\setbeamercolor{structure}{fg=CIDEBlue}
\setbeamercolor{alerted text}{fg=CIDEBlue}
\setbeamercolor{title}{bg=CIDEBlue,fg=white}
\setbeamercolor{subtitle}{fg=CIDEInk}
\setbeamercolor{frametitle}{bg=CIDEBlue,fg=white}
\setbeamercolor{palette primary}{bg=CIDEBlue,fg=white}
\setbeamercolor{palette secondary}{bg=CIDEBlue,fg=white}
\setbeamercolor{palette tertiary}{bg=CIDENavy,fg=white}
\setbeamercolor{palette quaternary}{bg=CIDENavy,fg=white}
\setbeamercolor{section in head/foot}{bg=CIDENavy,fg=white}
\setbeamercolor{section in head/foot shaded}{bg=CIDENavy,fg=white!50}
\setbeamercolor{mini frame}{bg=white,fg=white}
\setbeamercolor{mini frame shaded}{bg=white!35,fg=white!35}
\setbeamercolor{upper separation line head}{bg=CIDENavy}
\setbeamercolor{lower separation line head}{bg=CIDEBlue}
\setbeamercolor{block title}{bg=CIDEBlue,fg=white}
\setbeamercolor{block body}{bg=CIDEGrayLight,fg=CIDEInk}
\setbeamercolor{block title alerted}{bg=CIDENavy,fg=white}
\setbeamercolor{block body alerted}{bg=CIDEGrayLight,fg=CIDEInk}
\setbeamercolor{caption name}{fg=CIDEBlue}
\setbeamercolor{section in toc}{fg=CIDEBlue}
\setbeamercolor{section in toc shaded}{fg=CIDEGray}

% --- Tipografia ---------------------------------------------------------------
\setbeamerfont{author in head/foot}{size=\tiny}
\setbeamerfont{institute in head/foot}{size=\tiny}
\setbeamerfont{title in head/foot}{size=\tiny}

% --- Geometria general --------------------------------------------------------
\setbeamersize{text margin left=0.70cm,text margin right=0.70cm}
\setlength{\leftmargini}{1.2em}
\setlength{\leftmarginii}{1.1em}
\setlength{\abovecaptionskip}{0.25em}
\setlength{\belowcaptionskip}{0pt}

% Navegacion superior con nombres de seccion y mini-indicadores consecutivos.
\useoutertheme[subsection=false]{miniframes}
\setbeamertemplate{navigation symbols}{}
\setbeamertemplate{mini frame}[circle]
\setbeamertemplate{mini frame in current subsection}[circle]
% Colores explícitos: evita que algunas versiones de Beamer hereden el color
% del contenido de la diapositiva y vuelvan invisibles secciones futuras.
\setbeamertemplate{section in head/foot}{\color{white}\insertsectionhead}
\setbeamertemplate{section in head/foot shaded}{\color{white!55}\insertsectionhead}

% Barra superior determinista. Conserva el aspecto de miniframes y evita la
% oscilación de anchuras que presentan algunas versiones recientes de Beamer.
% Si se agregan o eliminan frames, sólo hay que actualizar estos rangos.
\newcount\cidenavdot
\newcommand{\cidenavdots}[2]{%
  \cidenavdot=#1\relax
  \loop
    \ifnum\value{framenumber}=\cidenavdot
      {\color{white}\raisebox{0.12ex}{\fontsize{4}{4}\selectfont$\bullet$}}%
    \else
      {\color{white!65}\raisebox{0.12ex}{\fontsize{4}{4}\selectfont$\circ$}}%
    \fi
    \kern-0.02em
    \ifnum\cidenavdot<#2\relax
      \advance\cidenavdot by 1\relax
  \repeat
}
\newcommand{\cidenavitem}[4]{%
  \parbox[c]{#1}{%
    \centering
    \ifnum\value{framenumber}<#2\relax
      \color{white!55}%
    \else
      \ifnum\value{framenumber}>#3\relax
        \color{white!55}%
      \else
        \color{white}%
      \fi
    \fi
    {\tiny\strut #4}\par\vspace{-0.55ex}%
    \cidenavdots{#2}{#3}%
  }%
}
\setbeamertemplate{headline}{%
  \leavevmode
  \begin{beamercolorbox}[wd=\paperwidth,ht=4.55ex,dp=0.95ex,leftskip=0.16cm,rightskip=0.16cm]{section in head/foot}
    \hbox to \dimexpr\paperwidth-0.32cm\relax{%
      \cidenavitem{1.62cm}{2}{5}{Punto de partida}\hfill
      \cidenavitem{2.22cm}{6}{16}{Equilibrio competitivo}\hfill
      \cidenavitem{1.18cm}{17}{22}{Optimalidad}\hfill
      \cidenavitem{1.16cm}{23}{29}{Gasto público}\hfill
      \cidenavitem{1.32cm}{30}{36}{Productividad}\hfill
      \cidenavitem{2.55cm}{37}{44}{Impuestos distorsionantes}\hfill
      \cidenavitem{1.35cm}{45}{48}{Bienes públicos}\hfill
      \cidenavitem{0.84cm}{49}{52}{Síntesis}%
    }%
  \end{beamercolorbox}%
}

\setbeamertemplate{caption}[numbered]
\setbeamertemplate{section in toc}[circle]
\setbeamertemplate{subsection in toc}[circle]
\setbeamertemplate{itemize item}{\raisebox{0.15ex}{\tiny$\blacktriangleright$}}
\setbeamertemplate{itemize subitem}{\raisebox{0.15ex}{\tiny$\bullet$}}
\setbeamertemplate{enumerate item}[circle]

% Banda de titulo de cada diapositiva.
\setbeamertemplate{frametitle}{%
  \nointerlineskip
  \begin{beamercolorbox}[
    wd=\paperwidth,
    ht=3.05ex,
    dp=1.05ex,
    leftskip=0.42cm,
    rightskip=0.42cm
  ]{frametitle}
    \usebeamerfont{frametitle}\insertframetitle
  \end{beamercolorbox}%
}

% Pie en dos niveles: autor/institucion y nombre corto/paginacion.
\setbeamertemplate{footline}{%
  \leavevmode
  \vbox{%
    \offinterlineskip
    \hbox{%
      \begin{beamercolorbox}[wd=.72\paperwidth,ht=2.05ex,dp=.55ex,leftskip=.35cm]{palette secondary}
        \usebeamerfont{author in head/foot}\insertshortauthor
      \end{beamercolorbox}%
      \begin{beamercolorbox}[wd=.28\paperwidth,ht=2.05ex,dp=.55ex,rightskip=.35cm plus1fil]{palette secondary}
        \hfill\usebeamerfont{institute in head/foot}\insertshortinstitute
      \end{beamercolorbox}%
    }%
    \hbox{%
      \begin{beamercolorbox}[wd=.86\paperwidth,ht=1.85ex,dp=.50ex,leftskip=.35cm]{palette tertiary}
        \usebeamerfont{title in head/foot}\insertshorttitle
      \end{beamercolorbox}%
      \begin{beamercolorbox}[wd=.14\paperwidth,ht=1.85ex,dp=.50ex,center]{palette tertiary}
        \usebeamerfont{title in head/foot}\insertframenumber{} / \inserttotalframenumber
      \end{beamercolorbox}%
    }%
  }%
  \vskip0pt
}

% Logo CIDE automatico en la esquina inferior derecha.
\newcommand{\RutaLogoCIDE}{Figuras/General/logocide.png}
\setbeamertemplate{background}{%
  \vbox to \paperheight{%
    \vfill
    \hbox to \paperwidth{%
      \hfill
      \includegraphics[height=0.58cm,keepaspectratio]{\RutaLogoCIDE}%
      \hspace{0.18cm}%
    }%
    \vspace{0.68cm}%
  }%
}

% Portada.
\setbeamertemplate{title page}{%
  \vspace*{0.35cm}
  \centering
  \begin{beamercolorbox}[wd=\textwidth,sep=0.24cm,center,rounded=false]{title}
    \usebeamerfont{title}\inserttitle\par
  \end{beamercolorbox}
  \vspace{0.25cm}
  {\usebeamerfont{subtitle}\usebeamercolor[fg]{subtitle}\insertsubtitle\par}
  \vfill
  {\usebeamerfont{author}\insertauthor\par}
  \vspace{0.35cm}
  {\usebeamerfont{institute}\insertinstitute\par}
  \vfill
  {\usebeamerfont{date}\insertdate\par}
  \vspace{0.25cm}
}

% --- Comandos reutilizables ---------------------------------------------------
\newcommand{\FigurasSesion}{Figuras/03_Equilibrio General Estático}
\newcommand{\figura}[2][0.65\textheight]{%
  \includegraphics[width=\linewidth,height=#1,keepaspectratio]{\FigurasSesion/#2}%
}
\newcommand{\cidehighlight}[1]{\textcolor{CIDEBlue}{\bfseries #1}}
\newcommand{\fuentefigura}[1]{%
  \vspace{0.03cm}\par{\raggedright\tiny\textit{Fuente:} #1\par}%
}
\newcommand{\notafuentes}[1]{\note{[Sources] #1}}

% Los enlaces externos son azules; los internos conservan el color de Beamer.
\hypersetup{
  colorlinks=true,
  urlcolor=blue,
  linkcolor=.,
  citecolor=.,
  filecolor=.
}

% No se generan diapositivas automaticas al cambiar de seccion.
\newif\ifsectiontoc
\sectiontocfalse
\AtBeginSection[]{%
  \ifsectiontoc
    \begin{frame}{Contenido}
      \tableofcontents[currentsection,hideothersubsections]
    \end{frame}
  \fi
}

% =============================================================================
% METADATOS
% =============================================================================
\title[Macroeconomía II]{Sesión 3 \\ Equilibrio general en una economía cerrada}
\subtitle{Un modelo macroeconómico de un periodo}
\author[LECO]{Arturo López}
\institute[CIDE]{Centro de Investigación y Docencia Económicas (CIDE)}
\date[Otoño 2026]{Licenciatura en Economía, Otoño 2026}

% =============================================================================
\begin{document}
% =============================================================================

\begin{frame}
  \titlepage
  \notafuentes{Williamson (2018), cap. 5.}
\end{frame}

% =============================================================================
\section{Punto de partida}
% =============================================================================

\begin{frame}{De las decisiones individuales al equilibrio}
  En el capítulo anterior obtuvimos dos relaciones:
  \[
    N^s=N^s(w;\pi-T),
    \qquad
    N^d=N^d(w;K,z).
  \]

  \begin{itemize}
    \setlength{\itemsep}{0.55em}
    \item El consumidor elige consumo y oferta de trabajo.
    \item La empresa elige producción y demanda de trabajo.
    \item Ambas decisiones se estudiaron tomando el salario real como dado.
  \end{itemize}

  \begin{alertblock}{El paso de esta sesión}
    Determinar conjuntamente el salario real y las cantidades que hacen compatibles todas las decisiones.
  \end{alertblock}
  \notafuentes{Williamson (2018), transición de los caps. 4 a 5.}
\end{frame}

\begin{frame}{¿Qué construiremos en esta sesión?}
  \begin{enumerate}
    \setlength{\itemsep}{0.55em}
    \item Un equilibrio competitivo para una economía cerrada y de un periodo.
    \item La frontera de posibilidades de producción de la economía.
    \item La equivalencia entre equilibrio competitivo y óptimo de Pareto.
    \item Los efectos del gasto público y de la productividad.
    \item Las distorsiones generadas por un impuesto proporcional al trabajo.
    \item Una extensión en la que el gobierno produce bienes públicos.
  \end{enumerate}
  \notafuentes{Williamson (2018), objetivos de aprendizaje del cap. 5.}
\end{frame}

\begin{frame}{Una economía cerrada y de un solo periodo}
  \begin{itemize}
    \setlength{\itemsep}{0.55em}
    \item \textbf{Economía cerrada:} no hay comercio ni flujos financieros con el resto del mundo.
    \item \textbf{Un periodo:} no existen ahorro, inversión ni decisiones intertemporales.
    \item \textbf{Bien único:} el producto puede destinarse a consumo privado o compras públicas.
    \item \textbf{Dos insumos:} capital \(K\), fijo en el periodo, y trabajo \(N\), determinado en equilibrio.
    \item \textbf{Competencia perfecta:} consumidor y empresa toman el salario real \(w\) como dado.
  \end{itemize}

  \vspace{0.2cm}
  \centering
  \cidehighlight{Es un laboratorio mínimo para estudiar mecanismos macroeconómicos.}
  \notafuentes{Williamson (2018), pp. 142--145.}
\end{frame}

\begin{frame}[t]{Un modelo transforma supuestos en resultados}
  \begin{columns}[T,onlytextwidth]
    \begin{column}{0.43\textwidth}
      \begin{itemize}
        \setlength{\itemsep}{0.5em}
        \item Variables exógenas: \(G\), \(z\) y \(K\).
        \item El modelo determina \(C\), \(Y\), \(N\), \(w\), \(T\) y \(\pi\).
        \item Un experimento modifica una variable exógena y sigue la respuesta del equilibrio.
      \end{itemize}
    \end{column}
    \begin{column}{0.53\textwidth}
      \centering
      \figura[4.0cm]{fig_5_1_modelo_exogenas_endogenas.png}
      \fuentefigura{Williamson (2018), figura 5.1.}
    \end{column}
  \end{columns}
  \notafuentes{Williamson (2018), figura 5.1 y pp. 144--146.}
\end{frame}

% =============================================================================
\section{Equilibrio competitivo}
% =============================================================================

\begin{frame}{Agentes y mercados}
  \begin{itemize}
    \setlength{\itemsep}{0.6em}
    \item \textbf{Consumidor representativo:} ofrece trabajo y demanda un bien de consumo.
    \item \textbf{Empresa representativa:} demanda trabajo y produce el bien.
    \item \textbf{Gobierno:} compra \(G\) unidades del bien y recauda impuestos \(T\).
  \end{itemize}

  \vspace{0.25cm}
  Dos mercados deben vaciarse simultáneamente:
  \[
    \underbrace{N^s=N^d}_{\text{mercado de trabajo}},
    \qquad
    \underbrace{Y = C+G}_{\text{mercado de bienes}}.
  \]
  \notafuentes{Williamson (2018), pp. 143--147.}
\end{frame}

\begin{frame}{El gobierno también enfrenta una restricción}
  El gobierno compra \(G\) unidades del bien y las financia con un impuesto de suma fija:
  \[
    \boxed{G=T}.
  \]

  \begin{itemize}
    \setlength{\itemsep}{0.6em}
    \item \(G\) es exógeno: la política fiscal se fija fuera del modelo.
    \item \(T\) se ajusta para satisfacer el presupuesto gubernamental.
    \item Un impuesto de suma fija reduce el ingreso, pero no altera las decisiones del consumidor (y por lo tanto es un impuesto no distorsionante).
  \end{itemize}
  \notafuentes{Williamson (2018), sección ``Government'', pp. 143--144.}
\end{frame}


\begin{frame}{Definición general de equilibrio competitivo}
  Un \textbf{equilibrio competitivo} o \textbf{equilibrio walrasiano} es una asignación de recursos y un vector de precios tales que, dadas las variables exógenas:

  \begin{enumerate}
    \setlength{\itemsep}{0.45em}
    \item cada agente optimiza su objetivo sujeto a sus restricciones (conducta racional), tomando los precios como dados;
    
    \item todos los mercados se vacían.
  \end{enumerate}
  Cuando se nos pide definir un equilibrio competitivo para una economía en particular, debemos traducir estas condiciones a las características de la economía en cuestión.
\end{frame}

\begin{frame}{Propiedades de un equilibrio competitivo}
  La idea de equilibrio competitivo es que se cumplan las siguientes propiedades:
  \begin{enumerate}
    \item Todos los agentes están, individualmente, tomando las mejores decisiones que pueden.
    \item Las decisiones de todos los agentes son compatibles entre sí. Es decir, el vector de precios de equilibrio coordina dichas decisiones y garantiza que los mercados se vacíen de forma simultánea.
  \end{enumerate}
\end{frame}

\begin{frame}{Equilibrio general \textit{en general}}
  \begin{itemize}
    \item Lo que definimos anteriormente como equilibrio competitivo es un caso particular de equilibrio general.
    \item \cidehighlight{Equilibrio general} es un concepto más amplio.
    \item El \textbf{equilibrio competitivo} es un modelo de equilibrio general en el que hay competencia perfecta: todos los agentes son tomadores de precios.
  \end{itemize}
\end{frame}

\begin{frame}{Definición de equilibrio competitivo para esta economía}
  Un equilibrio competitivo en esta economía es una asignación
  \[
    \{C^*,N^{s*},N^{d*},Y^*,w^*,T^*,\pi^*\}
  \]
  y un vector de precios $(w^*)$ tales que, dadas las variables exógenas \((G,z,K)\):
  \begin{enumerate}
    \setlength{\itemsep}{0.45em}
    \item el consumidor maximiza su utilidad sujeto a sus restricciones, dado \(w^*\);
    \item la empresa maximiza beneficios sujeta a su tecnología, dado \(w^*\);
    \item el gobierno satisface \(G=T^*\);
    \item los mercados de trabajo y bienes se vacían.
  \end{enumerate}

  \notafuentes{Williamson (2018), pp. 144--147.}
\end{frame}

\begin{frame}{Ley de Walras}
  \begin{itemize}
    \item Sean \(M\) mercados y \(M\) ecuaciones de equilibrio. 
    \item La \cidehighlight{ley de Walras} establece que, si \(M-1\) mercados se vacían, el último mercado también se vacía.
    \item Es decir, la ley de Walras establece que una de las ecuaciones de equilibrio es redundante.
    \item En nuestro modelo, \(M=2\). Por lo tanto, si se vacía un mercado, el otro mercado también se vacía.
  \end{itemize}
\end{frame}


\begin{frame}{Definición formal de la Ley de Walras}
  \textit{Definición.} Sea una economía con \(M\) mercados y sea \(z_m(\mathbf{p}) \equiv q_m^d(\mathbf{p})-q_m^s(\mathbf{p})\) el exceso de demanda en el mercado \(m\), donde \(\mathbf{p}\in\mathbb{R}_{++}^{M}\) es el vector de precios de la economía.

  Cuando se satisfacen las restricciones presupuestarias de todos los
  agentes, la \cidehighlight{ley de Walras} establece que, para cualquier vector de precios \(\mathbf{p}\),
  \[
    \sum_{m=1}^{M} p_m z_m(\mathbf{p})=0.
  \]

  \begin{itemize}
    \item Por lo tanto, si \(M-1\) mercados se vacían, el mercado restante
    también debe vaciarse.

    \item En nuestro modelo existen dos mercados: bienes y trabajo.
    Definiendo $z_Y \equiv C+G-Y $ y $z_N \equiv N^d-N^s$, y tomando el bien como numerario, la ley de Walras implica
    \[
      z_Y+wz_N=0.
    \]
  \end{itemize}
\end{frame}

\begin{frame}{Ley de Walras en esta economía}
  \begin{proposition}
    Si el mercado laboral se vacía, el mercado de bienes también se vacía.
  \end{proposition}
  
  \textit{Demostración.} Sea \(N^{s*}=N^{d*}=N\) el nivel de empleo de equilibrio del mercado laboral. A partir de la recta presupuestaria del consumidor en equilibrio se tiene que
  $$
  C^*=w^*N^{s*} + \pi^* - T^* = w^*N + \pi^* - T^*
  $$
  Como en equilibrio se cumple que
  $$
  \pi^* = Y^* - w^*N, \quad G=T^*,
  $$
  entonces
  $$
  C^* = w^*N + Y^* - w^*N - G = Y^* - G.
  $$
  Por lo tanto, si se vacía el mercado laboral, también se vacía el mercado de bienes:
  $$Y^* = C^* + G.$$
\end{frame}



\begin{frame}[t]{De la función de producción a la FPP}
  \begin{columns}[T,onlytextwidth]
    \begin{column}{0.46\textwidth}
      En equilibrio, \(N=h-\ell\). Por tanto,
      \[
        Y=zF(K,h-\ell).
      \]
      Usando \(C=Y-G\):
      \[
        \boxed{C=zF(K,h-\ell)-G}.
      \]

      Esta relación es la frontera de posibilidades de producción (FPP) entre consumo y ocio.
    \end{column}
    \begin{column}{0.50\textwidth}
      \centering
      \figura[5.4cm]{fig_5_2_funcion_y_fpp.png}
      \fuentefigura{Williamson (2018), figura 5.2.}
    \end{column}
  \end{columns}
  \notafuentes{Williamson (2018), figura 5.2 y pp. 147--149.}
\end{frame}

\begin{frame}{La pendiente de la FPP es el producto marginal del trabajo}
  Partiendo de
  \[
    C(\ell)=zF(K,h-\ell)-G,
  \]
  se obtiene
  \[
    \frac{dC}{d\ell}
    =-zF_N(K,h-\ell)
    =-\mathrm{MPN}.
  \]

  \begin{itemize}
    \setlength{\itemsep}{0.55em}
    \item Una hora adicional de ocio reduce el empleo en una hora.
    \item El costo tecnológico es la producción que deja de generarse: el MPN.
    \item Como el MPN aumenta cuando cae el empleo, la FPP se vuelve más inclinada conforme aumenta el ocio.
  \end{itemize}
  \notafuentes{Williamson (2018), pp. 147--150.}
\end{frame}


\begin{frame}[t]{Representación gráfica del equilibrio competitivo}
\begin{columns}[T,onlytextwidth]
\small
\begin{column}{0.43\textwidth}
En el punto \(J\):
\[
  \mathrm{MRS}_{\ell,C}(C^*,\ell^*)
  =
  w^*
  =
  \mathrm{MPN}(K,N^{d*}).
\]

\begin{itemize}
  \item La empresa elige \(N^{d*}\) de modo que
  $
    \mathrm{MPN}(K,N^{d*})=w^*.
  $
  \item El consumidor elige consumo y ocio de modo que
  $
    \mathrm{MRS}_{\ell,C}(C^*,\ell^*)=w^*.
  $
  \item El salario real \(w^*\) vacía el mercado de trabajo:
  $
    N^{d*}=N^{s*}=h-\ell^*.
  $
  \item En consecuencia, también se vacía el mercado de bienes:
  $
    Y^*=C^*+G.
  $
\end{itemize}
\end{column}

\begin{column}{0.53\textwidth}
  \centering
  \figura[5.35cm]{fig_5_3_equilibrio_competitivo.png}
  \fuentefigura{Williamson (2018), figura 5.3.}
\end{column}

\end{columns}

\notafuentes{Williamson (2018), figura 5.3 y pp. 149--151.}
\end{frame}

% =============================================================================
\section{Optimalidad}
% =============================================================================

\begin{frame}{Optimalidad}
  \begin{itemize}
    \item En la sección anterior caracterizamos el equilibrio competitivo de la economía.
    \item Ahora nos preguntamos si la asignación de recursos generada por el equilibrio competitivo es \textit{deseable}: para esto necesitamos algún criterio de evaluación.
    \item En economía, el principal criterio de eficiencia es el \cidehighlight{óptimo de Pareto}.
    \item Una asignación es \cidehighlight{eficiente en el sentido de Pareto} si no es posible mejorar la situación de un agente sin empeorar la de otro.
  \end{itemize}
\end{frame}

\begin{frame}{Definición de Óptimo de Pareto}
  \textit{Definición.} Sea una economía con \(I\) agentes $i \in \mathcal{I} = \{1, \ldots, I\}$, $L$ bienes $l \in \mathcal{L} = \{1, \ldots, L\}$ y sea \(u^i(\mathbf{x}^i)\) la función de utilidad del agente \(i\), donde \(\mathbf{x}^i \in \mathbb{R}^L_+\) es una asignación factible de bienes para el agente \(i\). Una asignación factible \(\{\mathbf{x}^{i*}\}_{i=1}^I\) es un \cidehighlight{óptimo de Pareto} si no existe otra asignación factible \(\{\mathbf{\hat x}^i\}_{i=1}^I\) tal que
  $$
  u^i(\mathbf{\hat x}^i) \geq u^i(\mathbf{x}^{i*}) \quad \forall i \in \mathcal{I} \quad \text{y} \quad u^j(\mathbf{\hat x}^j) > u^j(\mathbf{x}^{j*}) \quad \text{para algún } j \in \mathcal{I}.
  $$
\end{frame}

\begin{frame}{Planificador social benevolente}
  \begin{itemize}
    \setlength{\itemsep}{0.55em}

    \item ¿Cómo encontramos la asignación eficiente en el sentido
    de Pareto para esta economía?
    \item Utilizaremos el recurso teórico de un \cidehighlight{planificador social benevolente}.
    \item El planificador social es una autoridad que busca maximizar el bienestar social de la economía.
    \item En nuestra economía con consumidor representativo, el bienestar
    social está dado por su función de utilidad. En una economía con varios
    consumidores, el planificador maximizaría una función de bienestar
    social que pondera la utilidad de cada uno.

    \item El planificador conoce las preferencias y la tecnología de producción. No tiene ningún tipo de restricción más que las restricciones de recursos de la propia economía. Tiene el poder de redistribuir y elegir la asignación que maximiza la utilidad social.

    \item Es posible demostrar que la asignación elegida por el planificador es eficiente en el sentido de Pareto.
  \end{itemize}
\end{frame}

\begin{frame}{El problema del planificador}

  \vspace{0.2cm}
  El planificador social resuelve el siguiente problema:
  \[
    \max_{C,\ell}\ U(C,\ell)
    \quad\text{sujeto a}\quad
    C+G\leq zF(K,N), \ \ N \leq h-\ell.
  \]

  Para una solución interior y dado que $U_C > 0$ y $U_\ell > 0$, las C.P.O. son
  \[
  \begin{aligned}
    U_C(C,\ell)-\lambda &=0,\\
    U_\ell(C,\ell)-\lambda zF_N(K,h-\ell)&=0,\\
    zF(K,h-\ell)-C-G&=0.
  \end{aligned}
  \]
  
  Por lo tanto,
  \[
    \boxed{
      \mathrm{MRS}_{\ell,C}
      =
      \mathrm{MPN}
    }
  \]
\end{frame}

\begin{frame}[t]{Representación gráfica del óptimo de Pareto}
  \begin{columns}[T,onlytextwidth]
    \begin{column}{0.43\textwidth}
      \begin{itemize}
        \setlength{\itemsep}{0.55em}
        \item La asignación eficiente está en la tangencia con la FPP.
        \item Los puntos interiores a la FPP desperdician recursos.
        \item Sobre la frontera, el planificador elige la curva de indiferencia más alta alcanzable.
        \item En esta asignación, la MRS es igual a la tasa marginal de transformación, que es MPN.
      \end{itemize}
    \end{column}
    \begin{column}{0.53\textwidth}
      \centering
      \figura[5.35cm]{fig_5_4_optimo_pareto.png}
      \fuentefigura{Williamson (2018), figura 5.4.}
    \end{column}
  \end{columns}
  \notafuentes{Williamson (2018), figura 5.4 y pp. 151--153.}
\end{frame}

\begin{frame}{El planificador y el mercado asignan recursos de la misma manera}
  \begin{itemize}
  \item Esto parece redundante, pues hemos encontrado que la condición de tangencia del planificador es la misma que la condición de equilibrio competitivo.
  \item Esto no es una coincidencia. Las asignaciones que resultan de un equilibrio competitivo cumplen siempre con la condición de la asignación del planificador.
  \item El siguiente resultado demuestra que, \cidehighlight{en general}, la asignación de recursos que resulta de un equilibrio competitivo es eficiente en el sentido de Pareto.
  \end{itemize}
\end{frame}

\begin{frame}{Primer teorema fundamental del bienestar}
\begin{theorem}{Primer teorema fundamental del bienestar.}
  Sea $\left(\mathbf{p}^*, \{\mathbf{x}^{i*}\}_{i=1}^I\right)$ un equilibrio competitivo. Si las preferencias de todos los agentes son convexas y localmente no saciables, entonces la asignación $\{\mathbf{x}^{i*}\}_{i=1}^I$ es eficiente en el sentido de Pareto.
\end{theorem}  
\end{frame}

\begin{frame}{Primer teorema fundamental del bienestar}
\small
\textit{Demostración.} Procedemos por contradicción. Supongamos que la asignación $\{\mathbf{x}^{i*}\}_{i=1}^I$ no es eficiente en el sentido de Pareto. Entonces existe otra asignación factible $\{\mathbf{\hat x}^i\}_{i=1}^I$ tal que
$$
  u^i(\hat{\mathbf{x}}^i)
  \geq
  u^i(\mathbf{x}^{i*})
  \quad\forall i,
  \qquad
  u^j(\hat{\mathbf{x}}^j)
  >
  u^j(\mathbf{x}^{j*})
  \quad\text{para algún }j.
$$

Como \(\mathbf{x}^{i*}\) maximiza la utilidad del consumidor \(i\)
sujeto a su restricción presupuestaria, entonces
\[
  \mathbf{p}^*\cdot\mathbf{x}^{i*}
  =
  \mathbf{p}^*\cdot\mathbf{e}^i.
\]
La no saciedad local implica
\[
  \mathbf{p}^*\cdot\hat{\mathbf{x}}^i
  \geq
  \mathbf{p}^*\cdot\mathbf{e}^i
  \quad\forall i,
\]
con desigualdad estricta para el consumidor \(j\). Por lo tanto,
\[
  \mathbf{p}^*\cdot\sum_{i=1}^{I}\hat{\mathbf{x}}^i
  >
  \mathbf{p}^*\cdot\sum_{i=1}^{I}\mathbf{e}^i.
\]

Lo cual viola la factibilidad y llegamos a una contradicción. Por lo tanto,
\(\{\mathbf{x}^{i*}\}_{i=1}^{I}\) es eficiente en el sentido de Pareto.
\hfill\(\square\)

\end{frame}


\begin{frame}{Interpretación del PTF}
  \begin{itemize}
    \item ¿Cuál es la lógica económica detrás del primer teorema fundamental del bienestar? ¿Por qué las elecciones del planificador social son las mismas que las que produce el mercado?
    \item El planificador social elige una asignación para maximizar la utilidad sujeta a las posibilidades tecnológicas. En una economía competitiva, cada hogar maximiza su utilidad sujeta a los precios. Sin embargo, esos precios reflejan a su vez las posibilidades tecnológicas. Por lo tanto, el hogar también está, indirectamente, maximizando la utilidad sujeta a las posibilidades tecnológicas.
  \end{itemize}
\end{frame}


\begin{frame}{Segundo teorema fundamental del bienestar}

\begin{theorem}
Suponga que las preferencias son continuas, localmente no saciadas y
convexas, y que los conjuntos de consumo y de producción son convexos. Entonces, para toda asignación Pareto eficiente existe un vector de precios \(\mathbf{p}^*\neq\mathbf{0}\) y un sistema de transferencias de suma fija \(\{\tau_i\}_{i=1}^{I}\), con \(\sum_i\tau_i=0\), tal que dicha asignación puede sostenerse con un equilibrio competitivo.
\end{theorem}

\begin{itemize}
  \item En nuestro modelo, las transferencias de suma fija desplazan la restricción presupuestaria sin modificar su pendiente \(-w^*\). Por ello, permiten redistribuir riqueza sin distorsionar las decisiones marginales.
\end{itemize}

\end{frame}

% =============================================================================
\section{Gasto público}
% =============================================================================

\begin{frame}{Política fiscal}
  Considere un aumento permanente de \(G_1\) a \(G_2\), financiado con impuestos de suma fija:
  \[
    \Delta T=\Delta G>0.
  \]

  La FPP cambia de
  \[
    C=zF(K,h-\ell)-G_1
  \]
  a
  \[
    C=zF(K,h-\ell)-G_2.
  \]

  \begin{itemize}
    \item Para cada nivel de ocio, el consumo privado factible cae exactamente en \(\Delta G\).
    \item La pendiente de la FPP no cambia para un mismo \(\ell\).
  \end{itemize}
  \notafuentes{Williamson (2018), pp. 157--160.}
\end{frame}

\begin{frame}[t]{Más gasto público genera un efecto ingreso negativo}
  \begin{columns}[T,onlytextwidth]
    \begin{column}{0.43\textwidth}
      \begin{itemize}
        \setlength{\itemsep}{0.48em}
        \item La FPP se desplaza hacia abajo.
        \item El consumidor reduce \(C\) y \(\ell\).
        \item Menos ocio implica más empleo.
        \item La producción aumenta, pero el salario real cae.
      \end{itemize}
    \end{column}
    \begin{column}{0.53\textwidth}
      \centering
      \figura[5.45cm]{fig_5_6_gasto_publico.png}
      \fuentefigura{Williamson (2018), figura 5.6.}
    \end{column}
  \end{columns}
  \notafuentes{Williamson (2018), figura 5.6 y pp. 157--160.}
\end{frame}

\begin{frame}{Predicciones de un aumento en \(G\)}
  \begin{center}
  \begin{tabular}{@{}lll@{}}
    \toprule
    \parbox[t]{3.2cm}{\raggedright\bfseries Variable} &
    \parbox[t]{2.8cm}{\raggedright\bfseries Respuesta} &
    \parbox[t]{7.1cm}{\raggedright\bfseries Mecanismo}\\
    \midrule
    \parbox[t]{3.2cm}{\raggedright Consumo \(C\)} &
    \parbox[t]{2.8cm}{\raggedright Disminuye} &
    \parbox[t]{7.1cm}{\raggedright Menor riqueza disponible para el sector privado.}\\[0.15cm]
    \parbox[t]{3.2cm}{\raggedright Empleo \(N\)} &
    \parbox[t]{2.8cm}{\raggedright Aumenta} &
    \parbox[t]{7.1cm}{\raggedright El consumidor reduce ocio para sostener consumo.}\\[0.15cm]
    \parbox[t]{3.2cm}{\raggedright Producción \(Y\)} &
    \parbox[t]{2.8cm}{\raggedright Aumenta} &
    \parbox[t]{7.1cm}{\raggedright Más trabajo eleva la producción.}\\[0.15cm]
    \parbox[t]{3.2cm}{\raggedright Salario \(w\)} &
    \parbox[t]{2.8cm}{\raggedright Disminuye} &
    \parbox[t]{7.1cm}{\raggedright Con más empleo, cae el MPN.}\\[0.15cm]
    \bottomrule
  \end{tabular}
  \end{center}
  \notafuentes{Williamson (2018), pp. 157--160.}
\end{frame}

\begin{frame}{Crowding out y multiplicador del gasto}
  La identidad de recursos implica
  \[
    \Delta C=\Delta Y-\Delta G.
  \]

  Como el modelo predice \(\Delta C<0\) y \(\Delta Y>0\), entonces
  \[
    0<\Delta Y<\Delta G.
  \]

  Por tanto, para un cambio marginal,
  \[
    \boxed{0<\frac{dY}{dG}<1}.
  \]

  \begin{itemize}
    \item El gasto público desplaza parcialmente al consumo privado.
    \item El aumento en la producción evita un desplazamiento uno a uno.
  \end{itemize}
  \notafuentes{Williamson (2018), pp. 158--160.}
\end{frame}

\begin{frame}[t]{La Segunda Guerra Mundial como experimento fiscal}
  \begin{columns}[T,onlytextwidth]
    \begin{column}{0.38\textwidth}
      \begin{itemize}
        \setlength{\itemsep}{0.45em}
        \item Las compras públicas aumentaron de forma extraordinaria.
        \item La producción se elevó con fuerza.
        \item El consumo privado cayó sólo moderadamente.
        \item El patrón cualitativo coincide con el modelo.
      \end{itemize}
    \end{column}
    \begin{column}{0.58\textwidth}
      \centering
      \figura[5.25cm]{fig_5_7_segunda_guerra.png}
      \fuentefigura{Williamson (2018), figura 5.7.}
    \end{column}
  \end{columns}
  \notafuentes{Williamson (2018), figura 5.7 y pp. 160--161.}
\end{frame}


\begin{frame}[t]{Compras públicas y egresos totales no son lo mismo}
  \begin{columns}[T,onlytextwidth]
    \begin{column}{0.49\textwidth}
      \centering
      \textbf{Compras de bienes y servicios}

      \vspace{0.08cm}
      \figura[4.3cm]{fig_5_12_gasto_publico_pib.png}
    \end{column}
    \begin{column}{0.49\textwidth}
      \centering
      \textbf{Egresos públicos totales}

      \vspace{0.08cm}
      \figura[4.3cm]{fig_5_13_egresos_publicos_pib.png}
    \end{column}
  \end{columns}

  \fuentefigura{Williamson (2018), figuras 5.12 y 5.13.}
  \notafuentes{Williamson (2018), figuras 5.12--5.13 y pp. 167--170.}
\end{frame}

% =============================================================================
\section{Productividad}
% =============================================================================

\begin{frame}{Experimento: aumenta la productividad total de los factores}
  Considere
  \[
    z_2>z_1.
  \]

  Para cualquier cantidad de trabajo,
  \[
    z_2F(K,N)>z_1F(K,N),
  \]
  y también aumenta el producto marginal del trabajo:
  \[
    z_2F_N(K,N)>z_1F_N(K,N).
  \]

  \begin{itemize}
    \item La economía puede producir más con los mismos insumos.
    \item Trabajar se vuelve más productivo y el salario real de equilibrio aumenta.
  \end{itemize}
  \notafuentes{Williamson (2018), pp. 159--165.}
\end{frame}

\begin{frame}[t]{Mayor productividad desplaza la tecnología}
  \begin{columns}[T,onlytextwidth]
    \begin{column}{0.43\textwidth}
      \begin{itemize}
        \setlength{\itemsep}{0.55em}
        \item La función de producción se desplaza hacia arriba.
        \item Para cada \(N\), aumenta \(Y\).
        \item Para cada \(N\), también aumenta el MPN.
        \item En términos de \(C\) y \(\ell\), la FPP se expande y se vuelve más inclinada.
      \end{itemize}
    \end{column}
    \begin{column}{0.53\textwidth}
      \centering
      \figura[5.0cm]{fig_5_8_productividad.png}
      \fuentefigura{Williamson (2018), figura 5.8.}
    \end{column}
  \end{columns}
  \notafuentes{Williamson (2018), figura 5.8 y pp. 161--163.}
\end{frame}

\begin{frame}[t]{El nuevo equilibrio tiene más consumo y producción}
  \begin{columns}[T,onlytextwidth]
    \begin{column}{0.42\textwidth}
      El equilibrio se mueve de \(F\) a \(H\).

      \begin{itemize}
        \setlength{\itemsep}{0.48em}
        \item \(C\) aumenta.
        \item \(Y=C+G\) aumenta.
        \item El bienestar aumenta.
        \item El salario real aumenta.
        \item El efecto sobre \(\ell\) y \(N\) es ambiguo.
      \end{itemize}
    \end{column}
    \begin{column}{0.54\textwidth}
      \centering
      \figura[5.35cm]{fig_5_9_productividad_equilibrio.png}
      \fuentefigura{Williamson (2018), figura 5.9.}
    \end{column}
  \end{columns}
  \notafuentes{Williamson (2018), figura 5.9 y pp. 162--164.}
\end{frame}

\begin{frame}{Predicciones de un aumento en \(z\)}
  \begin{center}
  \begin{tabular}{@{}lll@{}}
    \toprule
    \parbox[t]{3.0cm}{\raggedright\bfseries Variable} &
    \parbox[t]{3.0cm}{\raggedright\bfseries Respuesta} &
    \parbox[t]{7.2cm}{\raggedright\bfseries Razón}\\
    \midrule
    \parbox[t]{3.0cm}{\raggedright \(C,Y\)} &
    \parbox[t]{3.0cm}{\raggedright Aumentan} &
    \parbox[t]{7.2cm}{\raggedright Se expande el conjunto factible.}\\[0.15cm]
    \parbox[t]{3.0cm}{\raggedright Bienestar} &
    \parbox[t]{3.0cm}{\raggedright Aumenta} &
    \parbox[t]{7.2cm}{\raggedright Se alcanza una curva de indiferencia más alta.}\\[0.15cm]
    \parbox[t]{3.0cm}{\raggedright \(w\)} &
    \parbox[t]{3.0cm}{\raggedright Aumenta} &
    \parbox[t]{7.2cm}{\raggedright Aumenta la productividad marginal del trabajo.}\\[0.15cm]
    \parbox[t]{3.0cm}{\raggedright \(N\)} &
    \parbox[t]{3.0cm}{\raggedright Ambigua} &
    \parbox[t]{7.2cm}{\raggedright Se oponen los efectos ingreso y sustitución.}\\
    \bottomrule
  \end{tabular}
  \end{center}
  \notafuentes{Williamson (2018), pp. 162--165.}
\end{frame}

\begin{frame}[t]{Productividad: efectos sustitución e ingreso}
  \begin{columns}[T,onlytextwidth]
    \begin{column}{0.42\textwidth}
      \begin{itemize}
        \setlength{\itemsep}{0.42em}
        \item \(A\to D\): sustitución.
          \begin{itemize}
            \item el ocio se encarece;
            \item \(\ell\downarrow\), \(N\uparrow\).
          \end{itemize}
        \item \(D\to B\): ingreso.
          \begin{itemize}
            \item aumenta la riqueza;
            \item \(\ell\uparrow\), \(N\downarrow\).
          \end{itemize}
        \item Ambos efectos elevan \(C\).
      \end{itemize}
    \end{column}
    \begin{column}{0.54\textwidth}
      \centering
      \figura[5.35cm]{fig_5_10_productividad_efectos.png}
      \fuentefigura{Williamson (2018), figura 5.10.}
    \end{column}
  \end{columns}
  \notafuentes{Williamson (2018), figura 5.10 y pp. 163--165.}
\end{frame}

\begin{frame}[t]{El PIB y el residuo de Solow se mueven juntos}
  \begin{columns}[T,onlytextwidth]
    \begin{column}{0.39\textwidth}
      \begin{itemize}
        \setlength{\itemsep}{0.47em}
        \item El residuo de Solow es una medida contable de la PTF.
        \item Sus desviaciones respecto a tendencia comueven con el PIB real.
        \item Este patrón motiva modelos de ciclos reales impulsados por productividad.
      \end{itemize}
    \end{column}
    \begin{column}{0.57\textwidth}
      \centering
      \figura[5.35cm]{fig_5_11_pib_residuo_solow.png}
      \fuentefigura{Williamson (2018), figura 5.11.}
    \end{column}
  \end{columns}
  \notafuentes{Williamson (2018), figura 5.11 y pp. 165--166.}
\end{frame}

\begin{frame}{Interpretar la productividad con cautela}
  \begin{itemize}
    \setlength{\itemsep}{0.58em}
    \item El residuo de Solow se calcula como la parte del crecimiento que no explican los insumos medidos.
    \item Por ello puede capturar tecnología, pero también utilización del capital, errores de medición y variaciones no observadas en la calidad de los insumos.
    \item Que \(z\) y \(Y\) comuevan no prueba por sí mismo que los choques tecnológicos causen todas las fluctuaciones.
    \item La teoría produce una predicción conjunta para \(C,Y,N,w\) y bienestar; el contraste debe considerar todas esas variables.
  \end{itemize}

  \begin{alertblock}{Lección}
    Una variable residual no debe interpretarse mecánicamente como un choque tecnológico exógeno.
  \end{alertblock}
  \notafuentes{Williamson (2018), pp. 165--166 y discusión del residuo de Solow del cap. 4.}
\end{frame}

% =============================================================================
\section{Impuestos distorsionantes}
% =============================================================================

\begin{frame}{El impuesto de suma fija es un punto de referencia}
  Con un impuesto de suma fija, trabajar una hora adicional sigue generando \(w\) unidades de consumo:
  \[
    \mathrm{MRS}_{\ell,C}=w=\mathrm{MPN}.
  \]

  Con un impuesto proporcional \(t\) sobre el ingreso laboral, el consumidor recibe sólo
  \[
    w(1-t)
  \]
  por cada hora adicional de trabajo.

  \begin{alertblock}{Cuña fiscal}
    El impuesto separa el rendimiento social del trabajo, \(\mathrm{MPN}\), del rendimiento privado, \(w(1-t)\).
  \end{alertblock}
  \notafuentes{Williamson (2018), pp. 170--174.}
\end{frame}

\begin{frame}{El impuesto al trabajo modifica la decisión marginal}
  El consumidor enfrenta
  \[
    C=w(1-t)(h-\ell)+\pi.
  \]

  Su condición de optimalidad es
  \[
    \mathrm{MRS}_{\ell,C}=w(1-t).
  \]

  La empresa continúa fijando
  \[
    w=\mathrm{MPN}.
  \]

  Por tanto,
  \[
    \boxed{\mathrm{MRS}_{\ell,C}=(1-t)\mathrm{MPN}<\mathrm{MPN}}.
  \]

  \begin{itemize}
    \item El consumidor trabaja menos de lo socialmente eficiente.
  \end{itemize}
  \notafuentes{Williamson (2018), pp. 171--175.}
\end{frame}

\begin{frame}{Un modelo simplificado hace visible la distorsión}
  Suponga una tecnología lineal:
  \[
    Y=zN.
  \]

  Entonces
  \[
    \mathrm{MPN}=z,
    \qquad
    w=z,
    \qquad
    \pi=0.
  \]

  La restricción de recursos es
  \[
    C=z(h-\ell)-G,
  \]
  y el presupuesto del consumidor es
  \[
    C=z(1-t)(h-\ell).
  \]

  El presupuesto público exige
  \[
    \boxed{G=zt(h-\ell)}.
  \]
  \notafuentes{Williamson (2018), ecuaciones 5-7 a 5-12 y pp. 171--176.}
\end{frame}

\begin{frame}[t]{Tecnología lineal: FPP y demanda de trabajo}
  \begin{columns}[T,onlytextwidth]
    \begin{column}{0.49\textwidth}
      \centering
      \textbf{FPP lineal}

      \vspace{0.08cm}
      \figura[4.45cm]{fig_5_14_fpp_simplificado.png}
    \end{column}
    \begin{column}{0.49\textwidth}
      \centering
      \textbf{Demanda laboral perfectamente elástica}

      \vspace{0.08cm}
      \figura[4.45cm]{fig_5_15_demanda_trabajo_simplificado.png}
    \end{column}
  \end{columns}

  \fuentefigura{Williamson (2018), figuras 5.14 y 5.15.}
  \notafuentes{Williamson (2018), figuras 5.14--5.15 y pp. 171--173.}
\end{frame}

\begin{frame}[t]{El equilibrio con impuesto no es eficiente}
  \begin{columns}[T,onlytextwidth]
    \begin{column}{0.42\textwidth}
      \begin{itemize}
        \setlength{\itemsep}{0.5em}
        \item \(AB\): FPP de la economía.
        \item \(DF\): presupuesto después de impuestos.
        \item \(H\): equilibrio competitivo.
        \item \(E\): óptimo de Pareto.
        \item En \(H\), consumo, empleo y producción son demasiado bajos.
      \end{itemize}
    \end{column}
    \begin{column}{0.54\textwidth}
      \centering
      \figura[5.35cm]{fig_5_16_impuesto_trabajo.png}
      \fuentefigura{Williamson (2018), figura 5.16.}
    \end{column}
  \end{columns}
  \notafuentes{Williamson (2018), figura 5.16 y pp. 173--175.}
\end{frame}

\begin{frame}[t]{La recaudación no aumenta indefinidamente con la tasa}
  \begin{columns}[T,onlytextwidth]
    \begin{column}{0.42\textwidth}
      \[
        R(t)=zt\,[h-\ell(t)].
      \]

      \begin{itemize}
        \setlength{\itemsep}{0.46em}
        \item \(t=0\): no existe recaudación.
        \item \(t=1\): desaparece el rendimiento privado de trabajar.
        \item \(t^*\): maximiza la recaudación.
        \item Para \(G<R(t^*)\), pueden existir dos tasas que financien el mismo gasto.
      \end{itemize}
    \end{column}
    \begin{column}{0.54\textwidth}
      \centering
      \figura[5.05cm]{fig_5_17_laffer.png}
      \fuentefigura{Williamson (2018), figura 5.17.}
    \end{column}
  \end{columns}
  \notafuentes{Williamson (2018), figura 5.17 y pp. 175--177.}
\end{frame}

\begin{frame}{La curva de Laffer no implica que reducir \(t\) siempre recaude más}
  \begin{itemize}
    \setlength{\itemsep}{0.62em}
    \item A la izquierda de \(t^*\), una tasa mayor eleva la recaudación.
    \item A la derecha de \(t^*\), una tasa mayor reduce tanto la base gravable que la recaudación cae.
    \item Saber de qué lado se encuentra una economía es una pregunta empírica.
    \item La forma y el máximo dependen de la elasticidad de la oferta laboral y de otros márgenes de respuesta.
  \end{itemize}

  \begin{block}{Punto central}
    La recaudación depende de la tasa y de la respuesta endógena de la base tributaria.
  \end{block}
  \notafuentes{Williamson (2018), pp. 175--177.}
\end{frame}

\begin{frame}[t]{Dos tasas pueden financiar el mismo gasto}
  \begin{columns}[T,onlytextwidth]
    \begin{column}{0.42\textwidth}
      Para \(t_1<t_2\):
      \begin{itemize}
        \setlength{\itemsep}{0.47em}
        \item ambas tasas recaudan \(G\);
        \item con \(t_1\), el rendimiento de trabajar es mayor;
        \item el equilibrio \(F\) tiene más \(C,N,Y\);
        \item el consumidor alcanza mayor bienestar que en \(H\).
      \end{itemize}
    \end{column}
    \begin{column}{0.54\textwidth}
      \centering
      \figura[5.35cm]{fig_5_18_dos_equilibrios.png}
      \fuentefigura{Williamson (2018), figura 5.18.}
    \end{column}
  \end{columns}
  \notafuentes{Williamson (2018), figura 5.18 y pp. 176--177.}
\end{frame}

% =============================================================================
\section{Bienes públicos}
% =============================================================================

\begin{frame}{Ahora el gasto público genera utilidad}
  Para separar la elección del tamaño del gobierno, suponga una dotación fija \(Y\) y preferencias
  \[
    U(C,G).
  \]

  El gobierno transforma impuestos \(T\) en bienes públicos con eficiencia \(q\):
  \[
    G=qT.
  \]

  Como el consumidor enfrenta
  \[
    C+T=Y,
  \]
  la FPP entre consumo privado y bienes públicos es
  \[
    \boxed{C+\frac{G}{q}=Y}.
  \]
  \notafuentes{Williamson (2018), sección ``A Model of Public Goods'', pp. 177--181.}
\end{frame}

\begin{frame}[t]{El tamaño óptimo del gobierno equilibra costos y beneficios}
  \begin{columns}[T,onlytextwidth]
    \begin{column}{0.42\textwidth}
      En una solución interior,
      \[
        \frac{U_G}{U_C}=\frac{1}{q}.
      \]

      \begin{itemize}
        \setlength{\itemsep}{0.45em}
        \item \(U_G/U_C\): beneficio marginal del bien público.
        \item \(1/q\): consumo privado sacrificado por una unidad de \(G\).
        \item \(G^*\) maximiza el bienestar, no la recaudación.
      \end{itemize}
    \end{column}
    \begin{column}{0.54\textwidth}
      \centering
      \figura[5.25cm]{fig_5_19_gasto_optimo.png}
      \fuentefigura{Williamson (2018), figura 5.19.}
    \end{column}
  \end{columns}
  \notafuentes{Williamson (2018), figura 5.19 y pp. 178--180.}
\end{frame}

\begin{frame}[t]{Una economía más rica demanda más bienes públicos}
  \begin{columns}[T,onlytextwidth]
    \begin{column}{0.42\textwidth}
      Si aumenta \(Y\):
      \begin{itemize}
        \setlength{\itemsep}{0.52em}
        \item la FPP se desplaza hacia afuera en forma paralela;
        \item si \(C\) y \(G\) son normales, ambos aumentan;
        \item la participación \(G/Y\) depende de la elasticidad ingreso de los bienes públicos.
      \end{itemize}
    \end{column}
    \begin{column}{0.54\textwidth}
      \centering
      \figura[5.25cm]{fig_5_20_aumento_pib.png}
      \fuentefigura{Williamson (2018), figura 5.20.}
    \end{column}
  \end{columns}
  \notafuentes{Williamson (2018), figura 5.20 y pp. 179--181.}
\end{frame}

\begin{frame}[t]{Un gobierno más eficiente puede ser más grande}
  \begin{columns}[T,onlytextwidth]
    \begin{column}{0.42\textwidth}
      Si \(q\) aumenta:
      \begin{itemize}
        \setlength{\itemsep}{0.46em}
        \item la FPP rota hacia afuera;
        \item producir \(G\) cuesta menos consumo privado;
        \item el efecto sustitución eleva \(G\) y reduce \(C\);
        \item el efecto ingreso eleva ambos;
        \item \(G\) aumenta y \(C\) puede subir o bajar.
      \end{itemize}
    \end{column}
    \begin{column}{0.54\textwidth}
      \centering
      \figura[5.25cm]{fig_5_21_eficiencia_gobierno.png}
      \fuentefigura{Williamson (2018), figura 5.21.}
    \end{column}
  \end{columns}
  \notafuentes{Williamson (2018), figura 5.21 y pp. 180--181.}
\end{frame}

% =============================================================================
\section{Síntesis}
% =============================================================================

\begin{frame}{Tres experimentos, tres mecanismos}
  \small
  \centering
  \begin{tabular}{@{}lllll@{}}
    \toprule
    \parbox[t]{2.4cm}{\raggedright\bfseries Cambio} &
    \parbox[t]{1.5cm}{\raggedright\bfseries \(C\)} &
    \parbox[t]{1.5cm}{\raggedright\bfseries \(Y\)} &
    \parbox[t]{1.5cm}{\raggedright\bfseries \(N\)} &
    \parbox[t]{4.5cm}{\raggedright\bfseries Mecanismo principal}\\
    \midrule
    \parbox[t]{2.4cm}{\raggedright \(G\uparrow\), impuesto de suma fija} &
    \parbox[t]{1.5cm}{\raggedright Baja} &
    \parbox[t]{1.5cm}{\raggedright Sube} &
    \parbox[t]{1.5cm}{\raggedright Sube} &
    \parbox[t]{4.5cm}{\raggedright Efecto ingreso negativo.}\\[0.2cm]
    \parbox[t]{2.4cm}{\raggedright \(z\uparrow\)} &
    \parbox[t]{1.5cm}{\raggedright Sube} &
    \parbox[t]{1.5cm}{\raggedright Sube} &
    \parbox[t]{1.5cm}{\raggedright Ambiguo} &
    \parbox[t]{4.5cm}{\raggedright Ingreso y sustitución opuestos sobre trabajo.}\\[0.2cm]
    \parbox[t]{2.4cm}{\raggedright \(t\uparrow\), modelo simplificado} &
    \parbox[t]{1.5cm}{\raggedright Baja} &
    \parbox[t]{1.5cm}{\raggedright Baja} &
    \parbox[t]{1.5cm}{\raggedright Baja} &
    \parbox[t]{4.5cm}{\raggedright Cuña entre rendimiento privado y social.}\\
    \bottomrule
  \end{tabular}
  \notafuentes{Síntesis de Williamson (2018), cap. 5.}
\end{frame}

\begin{frame}{Resultados que debemos conservar}
  \begin{enumerate}
    \setlength{\itemsep}{0.48em}
    \item El equilibrio general exige consistencia entre decisiones óptimas y vaciado de mercados.
    \item La FPP resume la tecnología y la restricción agregada de recursos.
    \item Con competencia perfecta e impuestos de suma fija, el equilibrio es eficiente.
    \item El gasto público no productivo desplaza consumo, pero aumenta empleo y producción.
    \item La productividad eleva consumo, producción, salario y bienestar; su efecto laboral es ambiguo.
    \item Los impuestos proporcionales crean una cuña y generan pérdidas de eficiencia.
  \end{enumerate}
  \notafuentes{Williamson (2018), resumen del cap. 5.}
\end{frame}

\begin{frame}{El modelo ya permite hacer macroeconomía}
  Comenzamos con decisiones individuales:
  \[
    \text{preferencias} + \text{tecnología} + \text{política}.
  \]

  El equilibrio las traduce en resultados agregados:
  \[
    (C^*,Y^*,N^*,w^*).
  \]

  \begin{itemize}
    \setlength{\itemsep}{0.6em}
    \item Los precios coordinan las decisiones descentralizadas.
  \end{itemize}
  \notafuentes{Williamson (2018), transición de los caps. 5 a 6.}
\end{frame}

\begin{frame}{Referencias de la sesión}
  \begin{itemize}
    \setlength{\itemsep}{0.75em}
    \item Williamson, Stephen D. (2018). \textit{Macroeconomics}, 6th ed. Pearson. Capítulo 5: ``A Closed-Economy One-Period Macroeconomic Model''.
  \end{itemize}
\end{frame}

\end{document}
